\documentclass{report}
\usepackage{amsmath,amssymb}
\setlength{\parindent}{0mm}
\setlength{\parskip}{1em}
\begin{document}
\begin{center}
\rule{6in}{1pt} \
{\large HyeongKae Park \\
{\bf Physics-Based Preconditioning and Accurate Time Integration for A Moment-based Scale-Bridging Algorithm}}

Theoretical Division \\ MS B216 \\ Los Alamos National Laboratory \\ Los Alamos \\ NM 87545
\\
{\tt hkpark@lanl.gov}\\
Dana Knoll\\
Chris Newman\\
Rick Rauenzahn\\
Jeff Densmore\\
Allan Wollaber\end{center}

Accurate modeling of neutral particle transport behavior is of great
importance in many science and engineering applications.
Many numerical transport algorithms suffer from slow convergence in
optically thick regions, where particles undergo a large number of
scattering and/or absorption re-emission events. We are developing a
moment-based, scale-bridging algorithm that can mitigate slow convergence
of transport algorithms. Our algorithm utilizes the low-order (LO)
continuum equations that are consistently derived from the high-order
(HO) transport equation. Due to discrete consistency, the LO system can
be used not only for accelerating the HO solution, but also coupling to
other physics.

In this talk, we present physics-based preconditioning strategy that
utilizes the nonlinear elimination technique for the solution of the LO
problem. The nonlinear elimination technique can reduce coupled PDEs/ODE
to a scalar parabolic PDE via block Gauss elimination. This technique
allows one to eliminate stiff nonlinear coupling from the system and
hence reduces the number of required Newton iterations, with a trade-off
of more complex Jacobian and nonlinear residual evaluations. We can then
solve the LO system efficiently via multigrid preconditioned
Newton-Krylov method with physics-based preconditioning. We discuss
computational efficiency with different preconditioning operators.

We also discuss time stepping strategy of our moment-based scale-bridging
algorithm for time-dependent thermal radiative transfer problems. Our
goal is to obtain (asymptotically) second-order time accurate and
consistent solutions without nonlinear iterations between the HO and LO
systems within a time step.
We compare the accuracy and efficiency of a predictor-corrector algorithm
with the nonlinearly consistent algorithm.


\end{document}
